\section{合卷：合纵、连横与游士——言辞怎样走到山口之前}

函谷关外的道路不宽，能同时开来的车马、粮秣和军队却很多。战国外交的难处，正在于每一个使者都可以把“天下”的利害说得极大，而真正要出兵的国家仍须计算本国的城邑、仓廪、边境和继承。后世常以“合纵”“连横”概括这种竞争：前者大意是东方诸国联合以抗秦，后者多指秦与一国或数国分别结交、使其不入合纵。传世文字中又有“合从”“连衡”等写法，未必是一套当时人人遵守、边界分明的两派学说；它们更适合作为理解联盟聚散的后出名称。

这些名称不能把各国写成两排整齐的棋子。魏要守河洛与大梁，韩的都城和道路紧贴秦东出的门户，赵须顾太行与北方边地，齐能以东方财赋和海滨之利选择介入，楚虽有江汉纵深却难免于北面战线，燕则总要量度辽西、蓟与中原的距离。秦的关中腹地与函谷关使它能够择地防守，又让它特别在意近处的韩、魏城邑。盟约所能做的，是暂时改变这些计算的顺序；它不能取消山口、渡口和粮道。

\subsection{把许多恐惧写成一个盟约，仍须把军队带到同一处}

\begin{event}{公元前318年}{五国合纵攻秦}{可靠纪年}{函谷关及其东侧}\eventanchor{WQ-0318-HEZONG}{战国·五国合纵攻秦}
楚、赵、魏、韩、燕合兵向西，试图以共同出兵阻住秦的东进，终未能越过函谷关取得决定性战果。各国忧惧秦的压力，是结盟的结构性原因；但忧惧并不等于目标相同。对韩、魏而言，眼前是近邻城邑和通道的存亡；对赵、楚、燕而言，远道动员还要与本国边境和内部事务争夺粮饷、兵员与注意力。

函谷关把秦军的防御优势压缩到可以经营的关隘，也把东方联军的困难放大为协同问题：谁先到、谁供粮、谁愿意在别国最危险的方向承担损失，皆不能由盟书自动回答。直接后果是这次大规模合纵未能改变秦据关防守的格局；中期影响则是诸国更常在共同抗秦的名义与各自求缓、求地、求援的实际需要之间往返。它不是“纵约无用”的证据，而是联盟若无持续指挥、补给和共同风险分担，便难以把同一份恐惧变成同一支军队。
\end{event}

这次受挫也不等于秦只须闭关自守。秦若要由关中夺取东方城邑，仍得面对联军再起、齐楚牵制和近邻抵抗。前293年的\eventref{WQ-0293-YIQUE}{战国·伊阙之战}之所以沉重，正因伊、洛之间的峡口把韩魏的防线和秦东出的粮道压在同一处；一场胜利能否延续，取决于能否把守军、道路和新得城邑接进既有网络。合纵与连横并非脱离战争的口才，而是各国围绕这些网络不断改写出兵先后的办法。

\subsection{苏秦、张仪：传记保存的是传统，不是一份会谈录}

苏秦、张仪的名字，几乎同“合纵”“连横”不可分离。《史记·苏秦列传》《张仪列传》把二人的游历、说服与成败组织为彼此映照的人物故事；《战国策》又在不同国策中保存了大量以他们或同类说客为主角的长篇言辞。这些材料十分重要，因为它们让人看见战国人如何把利害、名分、恐惧和荣辱说成可供君主选择的语言；但它们主要是战国以后逐层汇集、整理的说辞与叙事，不能当作每次朝见逐字留下的记录。

所谓苏秦“佩六国相印”、张仪一席之间尽解纵约，正有传记文学所需要的整齐与戏剧性。较可把握的，是诸侯廷中确有跨国游说、馈赠、结交和人材争夺，游士也可能借一国的需求进入另一国的权力核心；难以由此坐实的，是某段辞令的原貌、某次转向只由一人造成，或人物履历中每一处前后相接的年月。把苏、张看作两条强大的政治记忆，比把他们看作能单手转动天下的操盘者，更接近材料的限度。

\begin{event}{公元前316年}{秦取巴、蜀}{可靠纪年}{巴蜀与关中西南通道}\eventanchor{WQ-0316-SHU-BA}{战国·秦灭巴蜀}
秦趁蜀国内乱，出兵兼并蜀、巴。传世叙事将张仪与司马错关于先伐韩、先取蜀的争论写得鲜明；它可以显示秦廷确实要在中原近攻与西南经营之间权衡，不能让我们逐句复原殿上的辩论。较稳固的事件核心是，秦以军事和行政扩张取得成都平原及西南通道，使关中以外的粮食、人口与物资有了新的来源。

这件事也校正了“连横”只靠缔约的想象。秦在外交上同远国暂时修好，是否真能向东持续施压，仍取决于可供运输、征发和驻守的腹地。巴蜀的取得不是后来所有胜利的单一原因，却是使能条件之一：它扩大了秦可调用的资源，也使秦廷在面对东方联盟时多了一层回旋余地。
\end{event}

\begin{event}{公元前312年}{丹阳、蓝田的秦楚战事}{可靠纪年}{丹阳、蓝田及秦楚之间}\eventanchor{WQ-0312-QIN-CHU}{战国·秦楚战争}
秦与楚先后在丹阳、蓝田交战，楚军受挫。后世叙事常把楚怀王是否因张仪许地而绝齐，写成一场由巧言骗成的外交悲剧；《史记·楚世家》《张仪列传》与《战国策》保存的说辞，正使这一情节格外完整。其时间、许地数目和对答的细部在不同文本中并不一致，不能视作当日会谈的速记。

楚的选择也不能只用“受骗”解释。它要在齐秦之间分配关系，又要守住西北方向和广阔国内的调度；秦则把外交施压同边境用兵并用。触发因素可以是齐秦关系的转折，决定战局的却还有秦楚接壤处的道路、军粮与动员。直接后果是楚在西部和汉水方向的压力加重；更长的影响，是楚虽仍为大国，却更难把一次结交或绝交化作稳定的安全。
\end{event}

这类故事里的游士，并非只卖弄辞采。他们携带各国的地理、军势、婚姻、宾客和朝廷人事的消息，在信息不完整时替君主排列风险；也会依附权贵、受制于馈赠，或因旧主失势而改投新国。商鞅由魏入秦、吴起往返魏楚、范雎自魏至秦，都说明人才流动可以带来制度与战场经验，而非只有外交口舌。只是外来者能否被任用，仍由君主、相权、宗族和军功集团的关系决定；游说不是悬在国家之上的自由职业。

\subsection{称号、会盟与离心：同盟为何总在成功之后变难}

\begin{event}{公元前288年}{秦齐互称西帝、东帝，旋即废号}{可靠纪年}{秦、齐及列国外交场域}\eventanchor{WQ-0288-DI}{战国·东西帝号}
秦昭襄王与齐湣王一度互称“西帝”“东帝”，不久又废去帝号。王号已在\eventref{WQ-0334-KINGS}{战国·齐魏称王}中被强国相互承认；更高的称号本可抬升外交位置，却也立刻触动其余诸国对单极秩序的戒心。废号不是礼仪无关紧要，恰说明名分必须有军力、盟友和可接受的利益分配相托，才能长期通行。

《史记·秦本纪》《田敬仲完世家》及《战国策》把这一段与苏秦晚年活动等叙事相连。可确认的是帝号一度出现又撤回；至于谁以哪一句话促成或破坏它，仍应放在后出说辞传统中阅读。人事转折不在于某位辩士神奇地改变了天下，而在于秦、齐的君主和大臣发现，夺取超越诸王的名号会让本可分别谈判的邻国更易结成共同的防备。
\end{event}

前284年燕、赵、秦、韩、魏共同攻齐的\eventref{WQ-0284-QI}{战国·燕破齐}，最能显示联盟的两面。各国并不需要对齐怀有同一份怨恨，仍能在齐孤立时一同进军；可是夺取城邑之后，燕军必须独自面对齐地的守城、征收与地方归附，联盟不能替它长期治理。到\eventref{WQ-0279-TIANDAN}{战国·田单复齐}，占领体系裂开，正说明共同破敌与共同守成是两种不同的政治技术。

秦昭襄王用范雎，后来以“远交近攻”概括对外次序，见\eventref{WQ-0266-FANJU}{秦·范雎入秦与远交近攻}。这句断语不该化作一把万能钥匙：远方国家未必总能安抚，近邻也不会因距离近便无力抵抗。它所触及的却是一个可由地理检验的限制——韩、魏、赵的城邑若被秦夺取，较容易同函谷关内外的道路、县邑与粮运相接；齐、楚即使一时交好，也不能替秦驻守太行或伊洛的缺口。外交的效力，因而总受占领与补给的距离裁定。

\begin{event}{公元前241年}{五国再度合纵攻秦}{可靠纪年}{函谷关与关东诸国}\eventanchor{WQ-0241-HEZONG}{战国·五国再度合纵}
赵、楚、魏、韩、燕再度合兵攻秦，仍未能突破函谷关。秦已置东郡并直接压迫中原，列国共同的危机感比前318年更为迫切；但赵须防北方与秦军，楚有南北广域，魏、韩的城邑首当其冲，燕距主要战场最远，各国无法轻易把本国最后的机动兵力交给统一指挥。

联军退却的直接后果，是各国又回到分别守土、分别求援的处境。中期影响并非秦因一场会战便注定统一，而是秦能继续利用对手无法同时互救的间隙，逐一施压。长平后的力量失衡、秦的郡县化前沿与列国各自的继承和财政难题，才共同构成日后统一战争的条件；把前241年写成合纵“最后一次失败”而略过其后选择，会把历史压得过于平直。
\end{event}

\begin{sourcenote}[把说辞放回它的文本和地面]
\sourcebook{史记·苏秦列传》《张仪列传》《范雎蔡泽列传》与诸国世家}提供人物、王年和外交转折的主要传世线索；\sourcebook{战国策}保留了最丰富的游说辞令；\sourcebook{资治通鉴}有助于把战争和会盟放回前后年代。三者皆非当时连续的外交档案，尤其《战国策》的国别编排、人物归属与长篇对答，须同《史记》、编年材料和现代史源研究互证。关隘、河流、太行通道、城邑距离和可见的战争后果，能帮助检验一种策略在何处可能有效；它们不能替我们证明某一句话曾在殿上原样说出。
\end{sourcenote}

\begin{turningpoint}[制度得失：能结盟，不等于能共同作战]
\term{所得}：游士和盟约使各国能越过本国朝廷的视野，较快识别强邻、交换信息、争取援军，并在力量不均时制造牵制。秦也可借远交、宾客与人才流动争取时间，把外部压力同内部动员衔接起来。

\term{代价}：说辞会把共同利害说得比实际更整齐，君主和大臣却仍要为本国的粮道、城邑、军功与继承负责。合纵一旦要求某国先出兵、久供粮或放弃近处利益，离心便有了具体的理由；连横一旦让一国暂时获利，也会加深他国的不信。战国外交最深的限制不在辞令是否华美，而在国家能力和地理位置能否承受它所许下的行动。
\end{turningpoint}
